\subsection{Discussão dos resultados}

Os mapas de calor adquiridos antes e depois do reposicionamento do ponto de acesso mostram que a qualidade da cobertura Wi-Fi em ambiente
residencial não depende apenas da potência do equipamento ou da presença de uma antena omnidirecional.
A análise dos resultados mostra que o padrão de radiação é fortemente dependente da arquitetura do imóvel, pela presença de paredes,
móveis e objetos metálicos e pela ocorrência de reflexões e cancelamentos por multipercurso, fenômenos explicados pela literatura de redes (TANENBAUM; FEAMSTER; WETHERALL, 2021).
Em outras palavras, a posição do roteador em uma residência não pode ser definida apenas pela conveniência de instalação,
mas sim por um diagnóstico do ambiente de propagação.
No mapa inicial, a cobertura apresentou uma distribuição irregular. Regiões próximas ao Access point têm níveis de sinal consideravelmente mais altos,
enquanto ambientes mais afastados e obstruídos tiveram atenuação marcada.
Esse comportamento vai ao encontro com o modelo teórico apresentado na fundamentação: a atenuação por distância e os efeitos de absorção 
e reflexão alteram de forma significativa a potência recebida. A presença de obstáculos físicos e de superfícies refletoras gera regiões com queda acentuada do sinal, 
o que explica o motivo de uma rede aparentemente bem localizada pode ter pontos de baixa mesmo em áreas próximas ao roteador.
A discussão também sugere que a otimização da cobertura depende de uma análise do layout físico do imóvel. 
O ponto de acesso original foi instalado em um local de conveniência, mas não necessariamente no local que maximizava a penetração do sinal para os ambientes mais exigentes.
